\documentclass[10pt]{article}

\usepackage[utf8]{inputenc}

\usepackage[T1]{fontenc}

\usepackage{amsmath}

\usepackage{amsfonts}

\usepackage{amssymb}

\usepackage[version=4]{mhchem}

\usepackage{stmaryrd}

\usepackage{graphicx}

\usepackage[export]{adjustbox}

\graphicspath{ {./images/} }



\begin{document}

формула $A$ разрешима на системе множеств $F_{1}, \ldots, F_{n}$. Если формула $A\left(p_{1}, \ldots, p_{n}\right)$ разрешима на любой системе непустых конечных множеств $F_{1}, \ldots, F_{n}$, то она наз. в с е гда раз p ө шा и о й или финитно о б щ е з н а ч и м о й. Это определение имеет прозрачный смысл: Ф. о. формулы $A$ означает, что можно решить любую задачу $A\left(\mathfrak{H}_{1}, \ldots, \mathfrak{H}_{n}\right)$, зная лишь множества допустимых возможностей задач $\mathfrak{A}_{1}, \ldots, \mathfrak{A}_{n}$.



Множество $M L$ всех финитно общезначимых пропозициональных формул замкнуто относительно выводимости в интуиционистском исчислении высказываний и содержит все формулы, выводимые в әтом исчислении. Тем самым $M L$ является промежуточной (или суперинтуиционистской, суперконструктивной) логикой, наз. л о г и к о й М е д в ед е в а. Эта логика содержит формулы, не выводимые в иңтуиционистском исчислении высказываний (такова, напр., формула ( $7 a \supset$ $\supset(b \vee c)) \supset((\neg a \supset b) \vee(\rceil a \supset c)))$. Логика Медведева обладает свойством дизъюнктивности: если формула вида $A \vee B$ финитно общезначима, то по крайней мере одна из формул $A$ или $B$ финитно общезначима. Если пропозициональная формула не содержит какоголибо из логич. знаков $7, \vee, \supset$, то она финитно общезначима тогда и только тогда, когда выводима в интуиционистском исчислении высказываний. Все финитно общезначимые формулы выводимы в классич. исчислении высказываний; в то же время, напр., классически выводимая формула $a \vee 7 a$ не является финитно общезнатимой. Получены характөризации логики Медведева в терминах алгебраич. моделөй. Понятие Ф. о. различными способами может быть распространено на предикатные формулы.



\includegraphics[max width=\textwidth, center]{2023_07_09_cfbb636db98e4582af89g-1}

т. $142, N_{2}$ 5, c. $1015-18,1966$, т. $169, \mathrm{~N}_{1} 1$, c. $20-23$.



ФП. ФИНИТНАЯ ФУнКцИЯ - функция, определенная в нек-рой области пространства $E^{n}$ и имеющая принадлежащий к этой области комшактный носитель. Точнее, пусть функция $f(x)=f\left(x_{1}, \ldots, x_{n}\right)$ определена на области $\Omega \subset E^{n}$. Н о с и т е л е м $f$ наз. замыкание множества точек $x \in \Omega$, для к-рых $f(x)$ отлично от нуля $(f(x) \neq 0)$. Таким образом, можно еще сказать, что Ф. ф. в $\Omega$ есть такая определенная на $\Omega$ функция, тто ее носитель $\Lambda$ есть замкнутое ограниченное множество, отстоящее от границы $\Gamma$ области $\Omega$ ва расстоянии большем, чем $\delta>0$, где $\delta$ достаточно мало.



Обычно рассматривают вепрерывно дифференцируемые $k$ раз Ф. ф., где $k$-заданное натуральное число. Еще чаще рассматривают бесконечно дифференцируемые Ф. Ф. Функция



$\psi(x)= \begin{cases}e^{-1 /\left(|x|^{2}-1\right)}, & |x|<1, x^{2}=\sum_{j=1}^{n} x_{i}^{2}, \\ 0, & |x| \geqslant 1,\end{cases}$



может служить примером бесконечно дифференцируемой Ф. ф. в области $\Omega$, содержащей в себе шар $|x| \leqslant 1$.



Множество всех бесконечно дифференци руемых Ф. ф. в области $\Omega \subset E^{n}$ обозначают $D$. Над $D$ определяют линейные функционалы (обобщенные функции). При помощи функций $v \in D$ определяют обобщенные решения краевых задач.



В теоремах, посвященных задачам ва нахождение обобщенных решений, часто важно знать, является ли пространство $D$ плотным в нек-ром конкретно заданном пространстве функций. Известно, напр., что если граница $\Gamma$ ограниченной области $\Omega \subset E^{n}$ достаточно гладкая, то $D$ плотно в пространстве функций



$\stackrel{0}{W}_{p}^{r}(\Omega)=\left\{f: f \in W_{p}^{r}(\Omega),\left.\frac{\partial^{s} f}{\partial n^{s}}\right|_{\Gamma}=0, s=0,1, \ldots, r-1\right\}$, $(1 \leqslant p<\infty)$, т. е. в пространстве функций класса $W_{p}^{r}(\Omega)$ Соболева, равных нулю на $\Gamma$ вместе со своими нормалыными п роизводными до порядка $r-1$ включительно $(r=1,2, \ldots)$.



Лum.: [1] В л а д и м и р о в В. С , Обобщенные функции в математической физике, 2 изд., М., 1979, [2] С о б о л е в С J., Некоторые применения функционального анализа в математи-

ческой физике, Новосиб., 1962.

С. M. Никольский.



ФИНИТНО АПIРОКСИМИРУЕМАЯ ГРУППА групша, аппроксимируемая конечными группами. Пусть $G$ - группа, $\rho$ - отношение (иначе говоря, предикат) между әлементами и множествами элементов, определенное на $G$ и всех ее гомоморфных образах (напр., бинарное отношение равенства элементов, бинарное отношениө «әлемент $x$ входит в подгруппу $y$ », бинарное отношение сопряженности элементов и т. П.). Пусть $K-$ класс групі. Говорят, что группа $G$ а п п р о к с и м иру у т с я г р уппами из н $о \rho$, если для любых элементов и множеств әлементов из $G$, не находящихся в отношении $\rho$, существует такой гомоморфизм группы $G$ на группу из $K$, при к-ром образы атих әлөментов и множеств тоже не находятся в отношении $\rho$. Аппроксимируемость относительно равенства әлементов наз. просто аппроксимируемостью, Группа тогда и только тогда аппроксимируется группами класса $K$, когда она вкладывается в декартово произведениө групп из $K$. Финитная аппроксимируемость относительно $\rho$ обозначается $\Phi A \rho$; в частности, если $\rho$ пробегает предикаты равенства, сопряженности, вхождения в подгруппу, вхождения в конечно порожденную подгруппу и т. п., то получаются свойства (и классы) ФАР, ФАС, ФАВ, ФАВ и т. П. Из наличия атих свойств в группе вытекает разрешимость соответствующей алгоритмич. проблемы.



Лит.: [1] к а $\mathrm{p}$ г а п о л а в М. И., М е р з л я к ов Ю. И., Основы теории групा, 3 изд., М., 1982 . Ю. И. Мерзляков.



ФИНИТНО АППРОКСИМИРУЕМАЯ ПОЛУГРУППА, резид у ально коне чная пол уг p у п п a,- полугруппа, для любых двух различных элементов $a$ и $b$ к-рой существует такой ее гомоморфизм $\varphi$ в конечную полугруппу $S$, что $\varphi(a) \neq \varphi(b)$. Свойство полугруппы $S$ быть Ф. а. п. эквивалентно тому, что $S$ - подпрямое произведение конечных полугрупп. Финитная ашпроксимируөмость является одним из важных условий конечности (см. Полугруппа c условием конечности); оно тесно свявано с алгоритмическими проблемами: если $S$ - конечно определенная Ф. а. п., то в ней алгоритмически разрешима проблема равенства слов.



Ф. а. І. будут свободные полугрупшы, свободные коммутативные полугрупшы, свободные n-ступенно нильпотентные полугрупшы, свободные инверсные полугруппы (как алгебры с двумя операциями), полурешетки, конечно порожденные коммутативные полугруппы [1], конечно порожденные полугруппы матриц над нильпотентным или коммутативным кольцом, конечно порожденныө $n$-ступенно нильпотентные в смысле Мальцева (см. Нильпотентная полугруппа) регулярные полугруппы [4], см. также Финитно аппрокси.мируемая аруппа.



Прямое проивведение, свободное проивведение, ординальная сумма (см. Связка полуәрупn), 0-прямое объединение произвольного множества Ф. а. п. сами будут Ф. а. п. Другие конструкции, вообще говоря, нө сохраняют финитную аппроксимируемость. Идеальное расширение Ф. а. п. $S$ при помощи произвольной Ф. а. п. будет Ф. а. п., если, напр., $S$ р е д у к т и вн а я, т. е. любые два различных әлемента из $S$ индуцируют различные левые и различные правые внутренние сдвиги, в частности, если $S$ с сокращением или инверсная. Полурешетка нек-рого семейства редуктивных Ф. а. п. будет Ф. а. ІІ.



Если $S$ - Ф. а. п., то все ее максимальные подгруппы ф্রинитно апшроксимируемы. Для нек-рых типов





\end{document}